\section{The sparse-cones assertion}\label{sec:sparse-cones}

We retain the contradiction hypotheses, including the
lower-dimensional induction bounds and their proper-flat
consequences.  For $p\in\R^d$ and $y\ne p$, write $\pi_p(y)=[y-p]$.
A directional curve family is a finite family of curves in
$\PP^{d-1}$; membership means membership in its union, and its
total degree is the sum of the component degrees.

\begin{theorem}[Sparse cones]\label{sc:theorem}
Along a sequence with $M=o(A)$, there cannot be families $C_p$,
one for each $p\in P$, of irreducible projective curves, each
spanning projective dimension at most two, such that
\[
 \deg C_p\le E=\delta N/A,\qquad \delta\to0,
\]
and $1-o(1)$ of the unordered pairs $\{p,q\}$ satisfy
$\pi_p(q)\in C_p$ or $\pi_q(p)\in C_q$.
\end{theorem}

Suppose these data exist. Apply the results of Section~\ref{sec:profiles}.
The positive-mass core alternative is already contradictory.
We must treat a common window and two widely separated windows.
The selected covered incidences retain the global terminal deletion and
the light-line restriction in Proposition~\ref{prof:pair-selection}.
For the common-window rank estimate we will also use auxiliary graphs
of arbitrary ordered pairs $(p,q)$ with $p\ne q$ on the same selected point sets.
That estimate will be proved without a coverage or light-line hypothesis.

We use affine charts of the label map
\begin{equation}\label{sc:label-map}
 \psi(x,y)=(x,\pi_x(y)).
\end{equation}
For each finite family under discussion, these charts may be chosen to
contain all actual labels and to meet every component image.
The map and its pullbacks have bounded degree on these charts;
constant factors in testing degree fit inside a chosen open
plateau.  A \emph{dense graph} means at least a fixed positive
fraction of the ordered product of the two selected pieces.
Assertions about dense graphs always include arbitrary further
fixed positive-density restrictions and subsequences.

\subsection{Product images and rank in a common window}

Consider two pieces with width $s$ and the same limiting profile,
or the outermost pool allowed by Proposition~\ref{prof:pair-selection}.
Use $D_i,e_j,\sigma$ from~\eqref{prof:profile-functions}.
For an ordered graph $\Gamma$ on these pieces, write
$H(T)=H_{\psi(\Gamma)}(T)$ for the restriction rank of its labels.

We first make the passage from carrier pairs to product images precise.
The same construction will be used in the rank argument and, with one
fixed choice of parameters, at the outermost endpoint.  Its last estimate
applies to every later choice of image subfamily.

\begin{lemma}[Retained products and their image degree]
\label{sc:retained-products}
Fix a common positive-dimensional plateau after $j$ cuts, of dimension
$k=d-j<d-1$, and choose an interior time $\tau$ on that plateau.  Write
$n_x,n_y$ for the original source and target counts.  Given a fixed
$\varepsilon>0$, one may discard $o(n_x)$ source points and
$o(n_y)$ target points.  Write $m_x,m_y$ for the remaining counts, so
$m_x=(1-o(1))n_x$ and $m_y=(1-o(1))n_y$.  One may then discard carrier
pairs of total assigned product mass $o(m_xm_y)$, with the following conclusions.

Each remaining point is assigned to its unique smooth carrier in the
union of the two physical families.  Repeated carriers are pooled within
each family, and every flat carrier with positive assigned mass is real.
There are common physical equations for the whole original union whose Jacobian
has full codimension rank at these points.  A single polynomial
$g_{\rm ph}$, a linear combination of maximal Jacobian minors, is nonzero
at every retained point.  The equations and this test have degrees
$O(T_s(\tau))$.
When the two factors are the same pool, the point deletion and the
assignment may be taken to be the same in both roles.

Write $\mathcal R$ for the retained ordered carrier pairs and
\[
 Z_{X,Y}=\overline{\psi(X\times Y)}\qquad((X,Y)\in\mathcal R).
\]
Every $Z_{X,Y}$ has dimension $2k$, and at a general pair the connecting
direction is transverse to both physical tangent spaces.  Let $U$ be the
reduced union of the distinct $Z_{X,Y}$.  Then
\begin{equation}\label{sc:retained-image-upper}
 \deg U\le e_j^2\exp(o(s)).
\end{equation}
For any subfamily $\mathcal J$ of these distinct images, including one
chosen after further equations have been chosen, its represented carrier
pairs have normalized assigned mass at most
\begin{equation}\label{sc:retained-image-mass}
 \frac{1}{m_xm_y}
 \sum_{\substack{(X,Y)\in\mathcal R\\Z_{X,Y}\in\mathcal J}}n_Xn_Y
 \le \frac{\deg\bigl(\bigcup_{Z\in\mathcal J}Z\bigr)}{e_j^2}
       \exp(\varepsilon s+o(s)).
\end{equation}
Here $n_X,n_Y$ are the assigned masses after the point deletions.
\end{lemma}

\begin{proof}
Choose fixed times $\tau_{\rm eq}<\tau_{\rm test}<\tau$ between the start
of the plateau and $\tau$.  Use the earlier part of this interval for
the equation choices in Proposition~\ref{prof:regularity}, leaving the
strict gap up to $\tau_{\rm test}$ for its minor and proper-cut tests at
the sample points.  Apply its union version to the two physical carrier families.
It removes the negligible point sets excluded by its regularity conclusion
and supplies common equations for the whole union.  At each remaining
point there is now a unique smooth carrier.  Denote by $g_{\rm ph}$ the
common combination of maximal minors whose zeros were removed in this
step.  Discard every such point
whose carrier is a nonreal flat, using the degree-one proper sections in
the last paragraph of that proposition.  This costs another
$o(n_x)+o(n_y)$ points.  The source and target assignments are now
unambiguous.  Each assigned carrier belongs to the point's own original
piece path: that path already supplies a carrier through the point, and
uniqueness in the full union identifies it with the assigned member.
For two copies of the same pool take the union of the two negligible
point deletions in both roles; the resulting assignment is the same.
Extend the retained assignments by assigning each deleted point to a
containing carrier of its original piece before any later product tests,
using the same extension in both roles for a common pool.  Product testing
on the original sets then retains the present assignments under restriction, and its
$o(n_xn_y)$ pair loss is also $o(m_xm_y)$.  Pool equal retained carriers
within each of the two lists and omit members with zero mass.
All these choices have been made inside the interval ending at $\tau$,
which gives the stated degree bound.

The pooled budgets, with each original piece still supplying the bound
for its own points, give
\begin{equation}\label{sc:retained-carrier-weights}
 \begin{gathered}
 \max\left\{\sum_X\deg X,\sum_Y\deg Y\right\}\le e_j\exp(o(s)),\\
 \frac{n_X}{m_x}\le\frac{\deg X}{e_j}\exp(o(s)),\qquad
 \frac{n_Y}{m_y}\le\frac{\deg Y}{e_j}\exp(o(s)).
 \end{gathered}
\end{equation}
Indeed, the mass of a carrier in one piece is at most
$\deg(X)(N/e_j)\exp(o(s))$, and pooling $\exp(o(s))$ pieces changes
this by only $\exp(o(s))$.  The remaining point counts are still
$N\exp(-o(s))$.  A flat has degree one, so its number and its normalized
weight satisfy the hypotheses of Lemma~\ref{geo:generic-depth} with
$w=s$ and $e=e_j$.  The flat members are real by the preceding deletion.

Apply that lemma now with the fixed $\varepsilon$.  Let
$\mathcal Y_X$ be its target list for $X$, and set
$\mathcal R=\{(X,Y):Y\in\mathcal Y_X\}$.  The removed carrier pairs
have normalized product weight $o(1)$, hence account for $o(m_xm_y)$
actual pairs.  This deletion includes pairs generically tangent at either
endpoint.  Thus the target tangent projects injectively at a general pair,
and the first coordinate of $\psi$ records all source tangent directions;
the product image has dimension $2k$.  The preimages of a general
$2k$-fold linear section of that image are isolated away from the base
locus of the bounded-degree map on $X\times Y$.  The isolated-intersection
form of B\'ezout therefore gives
$\deg Z_{X,Y}\le C_d\deg X\deg Y$, and
\eqref{sc:retained-image-upper} follows from
\eqref{sc:retained-carrier-weights}.

We prove the uniform subfamily estimate while keeping these same lists.
Images belonging to different source carriers are distinct: the closure
of the first projection of $Z_{X,Y}$ is $X$.  Fix one source $X$ and one
of its distinct images $Z$.  To identify its source fibers, take its
closure $\widehat Z\subset X\times\PP^{d-1}$ before restricting the
direction coordinate to our affine chart.  Localizing the ideal of the
rational image at $\C(X)$ identifies the generic fiber with the closure
of the projection of $Y$ from the generic source point.  This fiber is
geometrically integral, since it is the image closure of the geometrically
integral variety $Y_{\C(X)}$.  After restricting $X$ to a dense open,
the projective family is flat with geometrically integral fibers of a
constant degree $d_Z$; this uses generic flatness, geometric integrality after
shrinking, and constancy of the Hilbert polynomial in a flat projective
family \cite[Tags~052A, 0559, 0C0E, and 0B9T]{stacksproject}.
We may also require that projection of $Y$ from each $x$ in this open set
has dimension $k$, by generic target transversality.
Its closed image is then a $k$-dimensional irreducible subvariety of the
$k$-dimensional integral fiber, so it equals the whole fiber.  Make these
restrictions simultaneously for the finite image and target lists.
Thus all targets with $Z_{X,Y}=Z$ have the same complete projected image
at a common general $x\in X$.  Write $t_x(Y)$ for their generic map degrees.
The fixed list bound says
\[
 \sum_{Y:\,(X,Y)\in\mathcal R,\ Z_{X,Y}=Z}t_x(Y)
 \le \exp(\varepsilon s).
\]
For a general $x$, the point-projection formula gives
$\deg Y-\operatorname{mult}_xY=t_x(Y)d_Z$.
The multiplicity is zero unless $X=Y$.  In that case $x$ is smooth and
the multiplicity is one; a coincident flat pair was deleted as tangential.
Consequently $\deg Y\le2t_x(Y)d_Z$ for every target in this sum.

Choose $k$ general affine hyperplanes in the source coordinates.  They
meet $X$ in $\deg X$ points of the preceding open set.  Then choose $k$
general hyperplanes in the direction coordinates; on each complete fiber
they give $d_Z$ points, with multiplicity.  The choices can be made so
all these points lie in the selected label chart and avoid intersections
between distinct images.  The resulting mixed section of $Z$ therefore
has $\deg X\,d_Z$ points.  The $2k$ blockwise hyperplanes are also affine
linear equations in the full label space.  Their homogenizations restrict
to linear equations on the ordinary projective closure of $Z$, where all
the counted affine points are isolated intersections.  The isolated-intersection
form of B\'ezout bounds their total multiplicity by the ordinary degree,
even if the same equations have additional components at infinity.
Applying this to a reduced image subfamily and summing the preceding
map-degree bound gives
\[
 \sum_{\substack{(X,Y)\in\mathcal R\\Z_{X,Y}\in\mathcal J}}
       \deg X\deg Y
 \le C_d\exp(\varepsilon s)
       \deg\bigl(\bigcup_{Z\in\mathcal J}Z\bigr).
\]
Multiplying the two weight bounds in
\eqref{sc:retained-carrier-weights} proves
\eqref{sc:retained-image-mass}.  No further generic-depth deletion is
made when $\mathcal J$ is chosen.
\end{proof}

\begin{lemma}[Common-window label rank]\label{sc:common-rank}
Let $\Gamma\subset\{(p,q)\in S_x\times S_y:p\ne q\}$ be any dense graph
on the selected pieces or pool; its pairs need not be covered by the $C_p$.
For its labels, at an interior testing time
before the last cut,
\begin{equation}\label{sc:label-rank}
 H(T_s(b))\ge
 \sigma(b)^2\min\left(1,\frac{T_s(b)}{D_1D_2}\right)
 \exp(-o(s)).
\end{equation}
The convention is a lower limit after dividing the logarithm
by $s$; in particular it permits any prescribed fixed positive
logarithmic tolerance, uniformly for dense subgraphs.
\end{lemma}

\begin{proof}
On the initial full-dimensional interval,
Lemma~\ref{prof:product-testing}, or its two successive
full-dimensional no-capture arguments, forces the testing-degree
envelope of a dense label graph to be the whole label space.
Its rank is comparable to $T^{2d-1}$.  Since $D_1D_2\ge T^2$
there, this is at least the right side of~\eqref{sc:label-rank}.

On a positive-dimensional plateau after $j\ge1$ cuts put
$k=d-j$.  Choose a physical regularity time between the start of this
plateau and the current testing time $b$.
When $k<d-1$, apply Lemma~\ref{sc:retained-products} with any fixed
$\varepsilon>0$.  When $k=d-1$, apply
Proposition~\ref{prof:regularity} to the union of the two physical
families, discard its rank-failure and nonreal-flat points, and assign each
survivor to its unique carrier from its original path.  Use the same
supported extension to the original point sets before product testing,
and pool equal retained carriers with their assigned masses.  The fixed
tangent-flat argument and the proper-flat
cap then discard carrier pairs generically tangent at either endpoint.
In both cases the point and carrier-pair losses remove only a negligible
fraction of the graph, and every remaining complete product image
$\overline{\psi(X\times Y)}$ has dimension $2k$.
Product testing forces these complete images into the testing-degree
envelope at almost all remaining graph pairs.

For this fixed envelope, choose one irreducible envelope component
containing each distinct complete image occurring at those surviving
pairs.  Let the pairs with that image inherit its choice.  Group all
envelope components by dimension and count each surviving pair in the
group of its chosen component; components receiving no pairs remain in
the degree of their group.  Every chosen component has dimension at
least $2k$, so some dimension $h\ge2k$ group carries a positive
fraction of the graph.  Write $a$ for its degree and $\Gamma'$ for
the graph portion it carries.

If $h\ge2k+1$, the Hilbert bounds give
\[
 H(T)\gtrsim aT^h,\qquad
 H_{\psi(\Gamma')}(T')\lesssim a(T'+1)^h.
\]
Borrow~\eqref{sc:label-rank} from a time just before the cut
starting this plateau.  The required right side grows with
exponent at most $2k+1$ within the plateau, so growth with
exponent $h$ suffices.  Endpoint approximation costs an
arbitrarily small fixed multiple of $s$ in the logarithm.

It remains to bound $a$ when $h=2k$.  Each image assigned to this group
is an irreducible $2k$-fold contained in its chosen irreducible
$2k$-dimensional envelope component, and therefore equals that component.
If $k=d-1$, necessarily $j=1$, and the projection of a generic
target hypersurface from a generic source point fills direction
space.  For each source carrier $X$ the product image is thus
$X\times\PP^{d-1}$ in these label coordinates.  A positive
mass of pairs requires a positive source mass.  The carrier
mass budget therefore gives
\[
 a\gtrsim e_1\exp(-o(s))=D_1\exp(-o(s)).
\]
The resulting rank $D_1T^{2d-2}\exp(-o(s))$ is sufficient:
on this plateau $T\le D_2$ to exponential accuracy and
\[
 D_1^2T^{2d-2}\min(1,T/(D_1D_2))
 \le D_1T^{2d-2}.
\]

Now suppose $k<d-1$.  The images assigned to this group represent
a positive amount of assigned product mass, since they carry
$\Gamma'$.  Apply \eqref{sc:retained-image-mass} to this subfamily.
Its degree is at most $a$, so
\[
 a\gtrsim e_j^2\exp(-\varepsilon s-o(s)).
\]
Since $\varepsilon$ is arbitrary, $a\ge e_j^2\exp(-o(s))$.
The Hilbert lower bound now gives $H(T)\ge e_j^2T^{2k}\exp(-o(s))$,
which is stronger than~\eqref{sc:label-rank} on this plateau.

These alternatives propagate through the finitely many distinct
cut times.  If the first boundary has degree $\exp(o(s))$ and
there is no earlier testing time of degree at least one, use
rank one there; the required boundary value is itself
$\exp(o(s))$.  This also handles several simultaneous first
cuts.  Every borrowed graph portion has fixed positive density
after passage to a subsequence, so the induction applies to it
as asserted.
\end{proof}

The graph covered in the selected orientation has an independent
upper bound, obtained by specializing at each actual center:
\begin{equation}\label{sc:common-upper}
 H(T)\le CNE(T+1).
\end{equation}
Near the last cut, $\sigma(b)$ is $N$ within an arbitrarily
small logarithmic tolerance.  Comparing
\eqref{sc:label-rank} and~\eqref{sc:common-upper}, the relevant
ratio is, to the same accuracy, at least
\begin{equation}\label{sc:common-ratio}
 \min\left(\frac{A}{\delta D_d},
            \frac{A}{\delta D_1D_2}\right).
\end{equation}
The second term has a strict positive $s$-scale margin because
$b_1+b_2<0$.  The first does also unless $s=\log B$ and $b_d=1$:
if the common window is not outermost, $s\le\tfrac12\log B$,
whereas in the outermost window its margin is $1-b_d$.
Choosing a fixed testing time sufficiently close to the final
cut therefore gives a contradiction in every other case.
No comparison between $\delta$ and an uncontrolled
$\exp(o(s))$ factor has been made.

\subsection{The outermost endpoint}

Assume now $s=\log B$ and $b_d=1$.  Use the common pool supplied by
Proposition~\ref{prof:pair-selection}; initially its size $n_0$ satisfies
\begin{equation}\label{sc:outer-pool}
 n_0\ge N\sqrt\delta,\qquad n_0=N\exp(-o(s)).
\end{equation}
Let $\Gamma_0$ be its graph covered in the selected orientation.  It has
at least $\alpha n_0^2$ pairs for a fixed $\alpha>0$, and retains the
terminal deletion and the light-line restriction.  Let $k$ be the
dimension of the final open plateau, $j=d-k$, and let $\beta=b_j<1$
be its starting time.  Balance gives
\begin{equation}\label{sc:outer-budgets}
 1\le k<d-1,\qquad e_j=B^{d-2k},\qquad N/e_j=A^k.
\end{equation}

At this endpoint the rank comparison can have zero logarithmic margin.
We will instead lift the covering direction curves to physical curves.
For this, we need one fixed family of product images and one common
smooth system containing the images used by all the lifts.

Choose fixed interior times
\begin{equation}\label{sc:outer-time-choice}
 \beta<\tau<t_0<t_1<t_2<1,
 \qquad t_0>\max\{\beta,1+b_1+b_2\}.
\end{equation}
Such choices exist because $b_1+b_2<0$.  Reserve the interval
$(\beta,\tau)$ for the finitely many physical regularity choices,
use $t_0$ for the image envelope and $t_1$ for the common label equations,
and reserve $(t_1,t_2)$ for the later pullback tests.  Now fix
\[
 0<\varepsilon<\tfrac14(t_1-t_0).
\]
Also fix the endpoint approximation and the logarithmic tolerances in
the rank argument below smaller than a quarter of both positive gaps
$t_0-\beta$ and $t_0-(1+b_1+b_2)$.

Apply Lemma~\ref{sc:retained-products} to the two copies of this pool,
with $\tau$ and $\varepsilon$ just fixed.  Use its same-pool deletion
and assignment, and retain its common physical equations and
$g_{\rm ph}$ for the full original union of final-plateau carriers.

When $k\ge2$, make one further point deletion before fixing the final
product family.  The carriers whose affine spans have dimension at most
$k+1$ number at most $e_j\exp(o(s))$, and each lies in a proper flat
containing at most $B^{k+1+o(1)}$ sample points.  The points assigned to
all these carriers therefore have total mass at most
\[
 B^{d-2k}B^{k+1+o(1)}=B^{d-k+1+o(1)}=o(n_0).
\]
Discard them in both roles.  No such deletion is made when $k=1$.

Write $S$ for the resulting pool, $n=|S|$, and $\mathcal C$ for its
positive-mass physical carriers.  Then
\begin{equation}\label{sc:outer-retained-pool}
 n=(1-o(1))n_0\ge(1-o(1))N\sqrt\delta,
 \qquad n=N\exp(-o(s)).
\end{equation}
Restrict the returned carrier pairs to $\mathcal C\times\mathcal C$
and call this fixed set $\mathcal R$.  Its omitted assigned product mass
is still $o(n^2)$.  The old assigned weights only decreased in the last
point deletion, while the normalizing point count changed by a factor
$1-o(1)$.  Thus \eqref{sc:retained-image-mass} remains valid, with the
same $\varepsilon$ and an adjusted $o(s)$, for every subfamily of the
distinct images
\[
 Z_{X,Y}=\overline{\psi(X\times Y)},\qquad (X,Y)\in\mathcal R.
\]
Let $U$ be their reduced union.  This is the final product-image family
for this endpoint; we will only take subfamilies of it.  In particular
$\deg U\le e_j^2\exp(o(s))$.  Restrict $\Gamma_0$ to $S\times S$ and
to the carrier pairs in $\mathcal R$, and call the result $\Gamma$.
The point and carrier-pair deletions together cost $o(n_0^2)$ pairs, so
$|\Gamma|\gtrsim n^2$.

When $k=1$, the unique carrier of each surviving target is its globally
assigned terminal curve from Proposition~\ref{prof:pair-selection}.
Indeed, the final plateau is then after the $(d-1)$st cut.  The curve
chosen in that global assignment belongs to the full original carrier
union used by the common physical equations, so uniqueness at the target
identifies it with the current carrier.  If that global assignment was
omitted because its chosen curve was a nonreal line, uniqueness would
make the current carrier that line, and the nonreal-flat point deletion
has removed the point.  Thus no globally unassigned target remains in
this case.

\begin{lemma}[Regular product images at the endpoint]
\label{sc:outer-regularity}
For the fixed $\mathcal R$ and $U$ above, there is a subfamily
$\mathcal R_{\rm sm}\subseteq\mathcal R$ whose omitted assigned product
mass is $o(n^2)$, and one system of label equations of degree at most
$B^{2-c}$ for some fixed $c>0$, with the following properties.
The equations contain every complete image $Z_{X,Y}$ for
$(X,Y)\in\mathcal R_{\rm sm}$.  Outside $o(n^2)$ further actual pairs
assigned to $\mathcal R_{\rm sm}$, their common zero locus is smooth of
dimension $2k$ at the label, the product map has differential rank
$2k$, and the connecting direction is transverse to both physical
tangent spaces.  If $k\ge2$, the two projected tangent spaces are also
unequal there.  The equations and every polynomial test used for these
actual-pair conclusions have degree at most $B^{2-c}$.
\end{lemma}

\begin{proof}
Put $T_i=T_s(t_i)$ and form
\[
 V=\operatorname{Env}_{T_0}(U).
\]
For each distinct image $Z\subset U$, choose once an irreducible component
$W(Z)$ of $V$ containing the whole image.  Every assigned pair represented
by $Z$ inherits that choice.  Since $\dim Z=2k$, a chosen component has
dimension at least $2k$; equality implies $W(Z)=Z$.

We first show that the images assigned to larger components represent
only $o(n^2)$ product mass.  Otherwise, after taking a subsequence, some
dimension $h\ge2k+1$ group represents a fixed positive normalized mass.
Form an auxiliary graph $\Lambda_h$ from all actual ordered pairs in
those assigned products, omitting the diagonal where $\psi$ is undefined.
Disjointness of the assignments makes its size equal to that product
mass before the omission, and the diagonal has at most $n$ pairs.
Thus $\Lambda_h$ is dense in $S\times S$.  This graph is independent of
the covered graph $\Gamma$; the arbitrary-graph scope of
Lemma~\ref{sc:common-rank} applies to it.
Every label of $\Lambda_h$ lies in the chosen dimension-$h$ component
containing its complete product image.

Write $a_h$ for the total degree of the dimension-$h$ component group of
$V$, including components receiving no images.  Its upper Hilbert bound
at an earlier degree $T'$ controls $H_{\psi(\Lambda_h)}(T')$, while its
lower envelope bound gives $H_V(T_0)\gtrsim a_hT_0^h$.  Borrow the
common-window label rank at a time just before $\beta$, with the fixed
endpoint tolerance chosen above.  If $\beta=-1$ and no earlier testing
degree is at least one, use the rank-one boundary convention in
Lemma~\ref{sc:common-rank}; then all $j$ earlier cut times equal $-1$
and the required boundary rank is one to exponential accuracy.
On the final plateau $\sigma(b)=e_jT_s(b)^k$.  Propagating with
$h\ge2k+1$ therefore gives
\[
 H_V(T_0)\ge e_j^2T_0^{2k}
 \min\left\{\frac{T_0}{T_s(\beta)},
             \frac{T_0}{D_1D_2}\right\}
       \exp(-\xi s-o(s)),
\]
where the fixed $\xi>0$ can be smaller than both gaps specified after
\eqref{sc:outer-time-choice}.  The two factors have normalized logarithms
$t_0-\beta$ and $t_0-(1+b_1+b_2)$, respectively, so the lower bound has
a strict positive margin.  But envelope equality and the upper Hilbert
bound on the pure $2k$-dimensional family $U$ give
\[
 H_V(T_0)=H_U(T_0)\le C e_j^2T_0^{2k}\exp(o(s)).
\]
This contradiction proves the claimed product-mass deletion.

Let $\mathcal I$ be the distinct images left after that deletion.
Each is an actual irreducible $2k$-dimensional component of the single
locus $V$, defined by equations of degree at most $T_0$; larger components
elsewhere in $V$ are kept in those equations.  Also
$\deg\mathcal I\le\deg U\le e_j^2\exp(o(s))$.
Apply Lemma~\ref{poly:regularity} to this pure component family at degree
$T_1$.  Its equations vanish on the whole family $\mathcal I$.
The exceptional image subfamily $\mathcal E$ has total degree at most
\[
 C\deg(\mathcal I)\frac{T_0}{T_1}
 \le e_j^2\exp(-(t_1-t_0)s+o(s)).
\]
Use \eqref{sc:retained-image-mass} for this particular subfamily of the
already fixed $U$.  Its represented normalized pair mass is at most
\[
 \exp(-(t_1-t_0-\varepsilon)s+o(s))=o(1).
\]
Here the same $\varepsilon$-lists chosen before $U$ are used; no new
generic-depth deletion depends on $\mathcal E$.  Let
$\mathcal R_{\rm sm}$ consist of the carrier pairs represented by
$\mathcal I\setminus\mathcal E$.  The two image deletions have omitted
only $o(n^2)$ assigned product mass.

The regularity lemma supplies a common polynomial $g$ of degree $O(T_1)$,
a linear combination of maximal Jacobian minors, which is nonzero on each
image in $\mathcal I\setminus\mathcal E$.  The single rational pullback
$g\circ\psi$ has a cleared numerator of degree $O(T_1)$ on the chosen
charts, nonidentically zero on every carrier pair in $\mathcal R_{\rm sm}$.
Product testing, with its
internal times in $(t_1,t_2)$, removes its actual zeros at a cost of
$o(n^2)$ pairs.  At every surviving label the common equations are locally
a smooth $2k$-fold.  They still vanish on all of $\mathcal I$, including
its exceptional components.

It remains to test the two tangent injectivities at actual pairs.
Write $f_a$ for the common physical equations chosen before $\tau$, and
$v=q-p$.  At a point $r\in S$, their Jacobian has kernel $T_r$, the tangent
space of its unique carrier.  Thus $v\notin T_r$ exactly when at least one
polynomial $\nabla f_a(r)\mathbin{\cdot}v$ is nonzero.  These tests have
degree $O(T_s(\tau))$ in $(p,q)$.  On every retained carrier pair the
source collection and the target collection each contain a generically
nonzero test, by the fixed generic-depth deletion.  For each endpoint,
choose a generic constant linear combination nonzero on every retained
product.  Product testing removes its actual zeros at cost $o(n^2)$.
The target injectivity, together with the source coordinate of $\psi$,
gives differential rank $2k$; the source injectivity is retained
separately for the secant step.

If $k\ge2$, every remaining physical carrier spans dimension at least
$k+2$ by the preceding carrier deletion.  Equality of the two projected tangent
spaces at a general retained pair would put both carriers in a common
$(k+1)$-flat by Lemma~\ref{geo:secants}, which is impossible.  This
inequality also has a bounded-degree test using the same physical
equations.  Put $c'=d-k$ and write $a_i(r)=\nabla f_i(r)$.  On the
full-rank transverse locus the row vectors
\[
 (a_i(r)\mathbin{\cdot}v)a_j(r)
 -(a_j(r)\mathbin{\cdot}v)a_i(r)
\]
span the $(c'-1)$-dimensional annihilator of $T_r+\mathbb Cv$.
Indeed, they lie in that annihilator, and choosing one row with nonzero
value on $v$ expresses every row combination annihilating $v$ in their
span.  The two projected tangent spaces are unequal exactly when the
combined row spans for $r=p,q$ have rank at least $c'$.  Their
$c'\times c'$ minors are polynomials of degree $O(T_s(\tau))$.  Choose a generic
combination nonzero on every retained product, and apply product testing
once more to remove its actual zeros.  All these tests fit in
the interval reserved below $t_2$.  Since $s=\log B$ and $t_2<1$, their
degrees and the common label-equation degrees are at most $B^{2-c}$ for
some fixed $c>0$.  This proves the lemma.
\end{proof}

Restrict $\Gamma$ by the carrier-pair and actual-pair deletions in
Lemma~\ref{sc:outer-regularity}.  It still has $\gtrsim n^2$ pairs.

Fix a center $p\in S$.  If a component $\gamma$ of $C_p$ is not contained
in the specialized common label equations, choose one such equation
nonzero on $\gamma$.  B\'ezout bounds its intersections with $\gamma$ by
$O(\deg\gamma\,B^{2-c})$.  Summing over these components gives
$O(EB^{2-c})$ directions.  Each supporting physical line has at most
$B^{o(1)}$ sample points, by the light-line restriction on $\Gamma$.
The corresponding target fraction is therefore at most
\[
 C\frac{EB^{2-c+o(1)}}{n}
 \le C\sqrt\delta\,B^{-c+o(1)}=o(1).
\]
This is uniform in $p$, so discarding all these incidences costs $o(n^2)$.

For a remaining component $\gamma\subset C_p$, take a marked pair
$(p,q)$ in the graph with $\pi_p(q)\in\gamma$.  Let $X,Y$ be its unique
physical carriers.  Near this label the common equations give a smooth
$2k$-fold $W$, and $\psi:X\times Y\to W$ has full differential rank at
$(p,q)$.  The analytic inverse function theorem therefore gives a local
biholomorphism.  On the affine open containing $q$ where $\pi_p$ is
regular, consider the algebraic preimage $Y\cap\pi_p^{-1}(\gamma)$.
Because $\psi$ records the source coordinate, its local inverse identifies
the reduced analytic germ of this preimage at $q$ with the germ of
$\gamma$ at $\pi_p(q)$.  That germ is pure one-dimensional.  Algebraic and
analytic local dimensions agree over $\C$, so an algebraic irreducible
component of the preimage containing any selected inverse branch is a
curve through $q$.  Let $C$
be its closure in $Y$.  The selected branch projects nonconstantly to
$\gamma$; hence $\overline{\pi_pC}=\gamma$, since $\gamma$ is irreducible.

If $k=1$, this curve is the target carrier $Y$ itself.  The identity
with the global terminal assignment established above now gives
$\overline{\pi_pY}=\gamma\subset C_p$, exactly an incidence removed by
Proposition~\ref{prof:pair-selection}.  This resolves $k=1$.

Suppose $k\ge2$.  For each remaining marked pair over $\gamma$, choose
one such irreducible lift, and pool repeated physical curves.  Retain its
carrier pair $(X,Y)\in\mathcal R_{\rm sm}$.  The physical minor test
$g_{\rm ph}(q)$ and the label and tangent tests are nonzero at its marked
point, so their restrictions are nonzero on the lift and hold on a
nonempty open part of it.  Choose one general direction on $\gamma$, common to this finite
family, avoiding the exceptional loci, branch values, and images of
intersections between distinct lifts.  Each lift then contributes its
generic map degree in distinct roots, and roots from distinct lifts are
also distinct.

There can be at most one such root in total.  At their common label the
same label equations are smooth of dimension $2k$ and contain every
complete product image attached to these lifts.  All products use the
same source carrier $X$: it is the unique member of the full physical
family at $p$.  If two roots $q',z'$ were distinct, the first product
would have full differential rank at $(p,q')$, and the complete second
product would also lie in the same local label system.  The
particular-pair assertion of Lemma~\ref{geo:secants} would then give
\[
 \overline{T_pX}\subseteq\overline{T_{q'}Y}
 \quad\text{in }\mathbb C^d/\mathbb C(q'-p).
\]
Both spaces have dimension $k$ by the two tangent-injectivity tests, so
they would be equal, contrary to the tangent-inequality test.  Hence
\begin{equation}\label{sc:lift-map-degree}
 \sum_{C\text{ distinct chosen lifts over }\gamma}
       \deg(C\longrightarrow\gamma)\le1.
\end{equation}

Finally, each lift loses at most half its physical degree under projection.
If $p\in C\subset Y$, uniqueness for the full original physical family
at $p$ forces $Y=X$.  The carrier $X$ is smooth at $p$, and the marked
point $q\in C$ has $q-p\notin T_pX$ by the source tangent test.
Corollary~\ref{poly:tangent-projection} therefore gives
$\operatorname{mult}_pC\le\deg C/2$.  If $p\notin C$, the multiplicity
is zero.  In both cases the point-projection formula yields
\[
 \deg\gamma\,\deg(C\to\gamma)
 =\deg C-\operatorname{mult}_pC\ge\tfrac12\deg C.
\]
Together with \eqref{sc:lift-map-degree}, this bounds the total degree of
the lifts over $\gamma$ by $2\deg\gamma$.  Every remaining target in the
row lies on a chosen lift.  The physical curve bound in
Lemma~\ref{poly:separators}, summed over $\gamma\subset C_p$, bounds their
number by $CEM$.  But
\[
 \frac{EM}{n}\le C\sqrt\delta\,\frac MA\longrightarrow0.
\]
Every remaining row thus has $o(n)$ partners, contradicting the surviving
$\gtrsim n^2$ graph.  This resolves the common outermost endpoint.

To treat separated windows, we next bound the projected degree at the
actual source centers.

\subsection{Projection depth for separated windows}\label{dep:section}

We retain the notation and contradiction regime \eqref{red:regime}.
Let $P_x,P_y$ be the two pieces selected by
Proposition~\ref{prof:pair-selection}, in separated windows of widths
$v$ and $s$, respectively. Thus $v,s\to\infty$, both are
$o(\log B)$, and either $v/s\to0$ or $v/s\to\infty$.
The limiting normalized drop profile is the same for both pieces.
Before using the target plateaux, omit from $P_y$ the points left
unassigned by the global terminal assignment in
Proposition~\ref{prof:pair-selection}. Its proof bounds their number
by $o(|P_y|)$ uniformly over the pieces, so this removes
$o(|P_x||P_y|)$ ordered pairs. Continue to write $P_y$ for the
retained target set. It has size $N\exp(-o(s))$ and inherits the
profile conclusions, since it is a $(1-o(1))$ fraction of the
selected piece. The proposition below uses this retained $P_y$.

For the target piece write
\[
 D_i=B\exp(b_i s),\quad e_j=\prod_{i\le j}D_i,
 \quad l=d-j,
\]
and consider a positive-dimensional plateau after at least one drop.
The path equations defining its current system have degrees
$B\exp(O(s))$. For this plateau, fix one containing
$l$-dimensional carrier for each point of $P_y$. On the final
one-dimensional plateau use its carrier from the global terminal
assignment. Pool repeated carriers and denote the resulting family
by $\mathcal Y$, with disjoint assigned target sets of masses $n_Y$.
Thus $\sum_{Y\in\mathcal Y}n_Y=|P_y|$. All later incidence and list
restrictions use this fixed assignment.
Our objective is a lower bound for the degree of the distinct
projected components carrying any fixed positive fraction of target
mass.  Several components may share a projected image, and each
projection may have degree greater than one.  We control this combined
multiplicity after deleting negligible incidence mass.

Give a center--component incidence $(p,Y)$ normalized weight
$n_Y/(|P_x||P_y|)$.  The weight of a collection is thus the probability
that a uniformly chosen center and a uniformly chosen assigned target
point determine an incidence in that collection.

We use the following precise interpretation of exponential accuracy.
An assertion with a factor $\exp(o(s))$ in a lower bound means that
for every fixed $\epsilon>0$ it holds with the smaller factor
$\exp(-\epsilon s)$ for all sufficiently large members of the
sequence, after the stated vanishing-mass deletions.
All finite choices of plateaux, density constants, and positive
logarithmic tolerances can be imposed simultaneously.

\begin{proposition}[Projected degree in separated windows]
\label{dep:projected-degree}
For each fixed positive tolerance $\epsilon$, one can discard
center--target-component incidences of total normalized weight
$o(1)$, and then discard $o(|P_x|)$ centers, so that the following
holds. For each remaining center $p$, every target component $Y$
whose incidence $(p,Y)$ remains satisfies
$\dim\overline{\pi_pY}=l$. For every subcollection of those components
carrying at least a fixed positive fraction of $|P_y|$, the distinct
projected components have summed degree at least
\begin{equation}\label{dep:degree-bound}
 \frac{e_j}{R}\exp(-\epsilon s).
\end{equation}
The assertion is uniform over those subcollections, including ones
chosen after $p$ and after subsequent equations have been chosen.
The following values of $R$ are valid:
\begin{align}
 R&=D_1 &&\text{in either orientation};
       \label{dep:first-R}\\
 R&=B\exp\left(-s\sum_{i>j}b_i\right)
       &&\text{if }v/s\longrightarrow0;
       \label{dep:narrow-R}\\
 R&=B\exp(-2s)
       &&\text{if }v/s\longrightarrow\infty,
          \quad b_{j+1}>0,\quad l<d-1.
       \label{dep:wide-R}
\end{align}
In \eqref{dep:degree-bound}, ``fixed positive fraction'' means any
constant $\rho>0$ independent of the sequence. After fixing
$\epsilon$ and any finite batch of plateaux and tolerances, the
incidence and center deletions are independent of the later fixed
choice of $\rho$; only the large-index threshold may depend on
$\rho$.
\end{proposition}

\begin{proof}
We first explain the degree and mass bookkeeping common to all three
cases, then prove their depth bounds.

\paragraph{Components, multiplicity, and incidence deletions.}
Proposition~\ref{prof:budget} gives
\begin{equation}\label{dep:carrier-budget}
 \sum_{Y\in\mathcal Y}\deg Y\le e_j\exp(o(s)),
 \qquad
 \sum_{Y\in\mathcal A} n_Y
 \le \deg(\mathcal A)\left(\prod_{i>j}D_i\right)\exp(o(s))
\end{equation}
for every subcollection $\mathcal A$.
Since $\sum_i b_i=0$ and $|P_y|=N\exp(-o(s))$, every subcollection
of fixed positive assigned mass has degree at least
$e_j\exp(-o(s))$.

Every $Y\in\mathcal Y$ is an irreducible component of the fixed target
path locus, by the component argument at the start of the proof of
Proposition~\ref{prof:regularity}. Larger components elsewhere remain
in the target equations.

Projection from $p$ can drop the dimension of $Y$ only if its generic
fiber is a line. For a fixed general smooth $y\in Y$, this puts $p$
in the proper affine tangent flat $y+T_yY$. The flat cap makes the
proportion of such centers at most $B^{-c_0+o(1)}$, uniformly for each
$Y$. Also, Corollary~\ref{poly:high-multiplicity} puts centers where
$\operatorname{mult}_pY>\deg(Y)/2$ in a hypersurface of bounded degree.
The initial full-dimensional no-capture property of the center piece
makes their proportion $o(1)$: its initial cutoff is $B^{1+o(1)}$,
well above every bounded degree. Average both conclusions with weights
$n_Y/|P_y|$. They remove only $o(1)$ of the center--component incidence
mass.

At every retained incidence, the point-projection formula gives
\begin{equation}\label{dep:projection-formula}
 t_p(Y)\deg\overline{\pi_pY}
       =\deg Y-\operatorname{mult}_pY
       \ge\tfrac12\deg Y,
\end{equation}
where $t_p(Y)$ is the generic map degree. For a fixed $p$, group
components having the same projected image $Z$ and define their depth
to be the sum of their map degrees. If all retained groups have depth
at most $R\exp(\eta s)$, then any subcollection of degree $a$ has
projected degree at least
\begin{equation}\label{dep:depth-to-degree}
 a\big/\bigl(2R\exp(\eta s)\bigr).
\end{equation}
Thus it suffices to bound depth after negligible weighted incidence
deletions. Fixed constants and the errors in
\eqref{dep:carrier-budget} are absorbed by choosing
$\eta>0$ sufficiently smaller than the desired $\epsilon$.

Whenever the deleted incidence mass is $\alpha=o(1)$, Markov's
inequality leaves all but $\sqrt\alpha$ of the centers with deleted
target mass at most $\sqrt\alpha$. Every later subcollection of
fixed positive mass still has positive mass after this deletion.
This is the uniformity mechanism in the proposition. No union bound
over the possible subcollections will be needed.

\paragraph{The first-cut bound.}
Use the system at the first plateau after a drop. Its equations have
degrees at most $D_1\exp(o(s))$, and every current descendant is
contained in this first system. The regularity conclusion of
Proposition~\ref{prof:regularity} makes its relevant carrier unique
and smooth at all but $o(|P_y|)$ assigned target points. For each such
point $q$, centers $p$ in its affine tangent flat form a
$B^{-c_0+o(1)}$ fraction of the source piece.

If, for a center and a current descendant $Y$, the generic joining ray
were contained in the first system, then the same containment holds
for every point of $Y$ where the ray is defined, by polynomial
identity. At an assigned regular point $q$ the line would lie locally
in its unique first carrier, so its direction would be tangent there.
The preceding averaging therefore discards all such
center--descendant incidences at negligible mass. On every other
generic ray at least one first-system equation restricts nontrivially.
Its degree bounds the number of distinct intersections with all
current target components by $D_1\exp(o(s))$. In characteristic zero,
the generic fibers defining the map degrees consist of distinct
points. This proves the depth bound for \eqref{dep:first-R}.
Only regularity of the carrier is used here; the original path
equations need not themselves have full Jacobian rank.

\paragraph{Centers in the narrower window.}
Suppose $v/s\to0$. The center piece can be simultaneously isolated as
zero-dimensional components of equations of degree
\[
 L=B\exp(o(s)),\qquad |P_x|=N\exp(o(s)),
\]
where the second error is signed. Fix $Y\in\mathcal Y$, and choose
$y$ general on it, also general for the projections from all finitely
many actual centers. For each center, count the additional moving
roots $z\ne y,p$ on target components whose images agree with
$\overline{\pi_pY}$. These points are generic on their components.
Here this includes avoiding their intersections with the other members
of the same finite family $\mathcal Y$. Their number is the group
depth minus the one fixed root $y$.

Parametrize such pairs by
\[
 p=y+t(z-y),\qquad z\in\bigcup_{Z\in\mathcal Y}Z.
\]
Impose the center-isolating equations on this parametrization.
At every counted solution the image point $p$ is isolated in the
center system. At its root $z$, the target union is locally just
the one member of $\mathcal Y$ containing that root, by the preceding
genericity condition. The generic ray has finite intersection with
that component, so the fiber is locally finite. Thus the solution is
isolated on the variety parametrized by $(z,t)$.
This variety has dimension $l+1$ and degree at most a dimensional
constant times $\sum_Z\deg Z$; the pulled-back equations have degree
at most $2L$. The isolated-intersection form of B\'ezout therefore
bounds the total number, over all centers, by
\begin{equation}\label{dep:narrow-root-count}
 C e_j\exp(o(s))L^{l+1}.
\end{equation}
Possible larger components of the center equations away from their
isolated sample points do not affect this count.

Divide \eqref{dep:narrow-root-count} by $|P_x|$. Its value, to the
stated accuracy, is
\[
 \frac{e_j B^{l+1}}{N}
 =B\exp\left(s\sum_{i\le j}b_i\right)
 =B\exp\left(-s\sum_{i>j}b_i\right)=R.
\]
The fixed root adds only one; $R=B^{1+o(1)}\to\infty$.
For every fixed logarithmic slack $\eta>0$, Markov's inequality bounds
the proportion of centers whose depth at this $Y$ exceeds
$R\exp(\eta s)$ by $\exp(-\eta s+o(s))$.
Weighted averaging in $Y$, followed by the row deletion already
described, proves \eqref{dep:narrow-R}.

\paragraph{Centers in the wider window: the generic lists.}
Suppose $v/s\to\infty$, $b_{j+1}>0$, and $l<d-1$.
Choose an interior center plateau just to the right of zero, before
the first strictly positive center drop. Its carrier dimension is
\[
 k=\#\{i:b_i>0\}\ge l,
 \qquad k<d.
\]
All componentwise proper cuts of degree $B\exp(Cs)$, with $C$ fixed,
are negligible for the center piece. To verify the degree comparison,
choose a fixed positive interior testing time $\tau$ on this center
plateau. Since $s=o(v)$,
$B\exp(Cs)=o(B\exp(\tau v))$, and leave a second interior time for
the pooling step of Proposition~\ref{prof:regularity}.

For each $p\in P_x$, fix one carrier $X(p)$ of this center plateau
which contains it. The distinct values form a source family
$\mathcal X$, with disjoint assigned sets
$S_X=\{p\in P_x:X(p)=X\}$. Apply the nonreal-flat deletion in
Proposition~\ref{prof:regularity} to this source assignment and to
the fixed target assignment on $\mathcal Y$. Use it to omit every
nonreal flat member from both input families, and omit any member
of zero assigned mass. The removed assignments have source mass
$r_x=o(|P_x|)$ and target mass $r_y=o(|P_y|)$. Discard every incidence whose center belongs to the
removed source set or whose target component is removed. Their
total normalized weight is at most
\[
 \frac{r_x}{|P_x|}+\frac{r_y}{|P_y|}=o(1).
\]
Call the surviving distinct families $\mathcal X_0,\mathcal Y_0$,
keeping their assigned sets, and put
\[
 n_x=\sum_{X\in\mathcal X_0}|S_X|=(1-o(1))|P_x|,
 \qquad
 n_y=\sum_{Y\in\mathcal Y_0}n_Y=(1-o(1))|P_y|.
\]
Both are at least $N\exp(-o(\log B))$, and every flat member of
these families is real. When $k=l$, the degree and assigned-mass
budgets, using $v,s=o(\log B)$, give at most $B^{d-l+o(1)}$ flat
members in each family and normalized weight at most
$B^{-(d-l)+o(1)}$ for each one. Thus
Lemma~\ref{geo:generic-depth} applies with error scale $\log B$ and,
when $k=l$, with $e=B^{d-l}$. Its removed product weight has
normalized incidence weight
\[
 \frac{n_xn_y}{|P_x||P_y|}\,o(1)=o(1)
\]
in the incidence normalization fixed above.

The lemma supplies a list $\mathcal Y_X\subseteq\mathcal Y_0$ for
each $X\in\mathcal X_0$ such that for every $Y\in\mathcal Y_X$,
a general line through $x\in X,y\in Y$ has
at most
\[
 q=B^{1/2}
\]
moving roots on that list. Using a smaller fixed exponent in the
lemma, if necessary, ensures this displayed bound for all sufficiently
large $B$. This is a bound on all possible generic list meetings,
not just on the roots visible in the finite sample. For a fixed
general $x$, a proper subvariety of an $l$-dimensional listed target
has projected dimension less than $l$, so cannot supply roots over
the generic direction image of $Y$.
For a retained actual center $p$, use only the list indexed by its
fixed carrier $X(p)$, further restricted by the earlier incidence
deletions at $p$.

We prove that excessive depth at an actual center is detected by a
proper cut on $X$ of degree $B\exp(O(s))$. Put
\[
 r=\left\lceil B\exp(-2s)\right\rceil.
\]
Because $s=o(\log B)$, we have $q=o(r)$ and $r=B^{1+o(1)}$.
Fix $X$ and $Y\in\mathcal Y_X$ for the following construction.

The generic list bound does not yet control actual centers.  We
encode intersections with a joining line by two binary forms.
Excessive retained depth will force a large rank decrease in their
Sylvester matrix.  Multiplicity compression will then place these
centers in a proper cut of degree fitting the center plateau.

\paragraph{A fixed generic target point and two binary forms.}
Write the current target path locus as
\[
 Z=V(f_1,\ldots,f_m),\qquad \deg f_i\le L_0=B\exp(O(s)).
\]
This is the fixed full path locus, including components omitted from
the generic-depth input family.
Every retained target is an irreducible component of $Z$, as proved
above. Choose $y\in Y$ general with two sets of conditions.
First, the generic list-root bound and the non-tangency conditions
must hold for generic $x\in X$. A dense open subset of $X\times Y$
has a dense open fiber over general $y$, so this is possible.
Second, for each of the finitely many retained actual centers $p$
assigned to $X$ and each listed target $Y'$ with $l$-dimensional image equal to
$\overline{\pi_pY}$, require all points of the generic fiber over
$\pi_py$ to be affine, distinct, different from $p$, and outside every
other irreducible component of $Z$.
The generically finite map from $Y'$ has a dense open image locus
with these properties: every excluded proper subset of $Y'$ has
image of dimension at most $l-1$. Pulling back these open sets to
$Y$ and intersecting the finite list preserves a nonempty open set.
Also avoid the finitely many actual centers. This fixes one $y$ for
all the specializations below.

Choose an integer $L\ge\max(L_0,B)$ with $L=B\exp(O(s))$.
Homogenize each path equation to its own degree and take all its
homogeneous monomial multiples to degree $L$. Their span has
projective base locus exactly $V(f_1^{\mathrm h},\ldots,f_m^{\mathrm h})$:
at every projective point some monomial of each prescribed degree is
nonzero. This base locus may contain extra components at infinity;
we retain all of them. The closure of each retained affine target
is still a component, since any larger component containing it
would meet the affine chart and contradict the component property.

Parametrize the line through $x$ and $y$ by
\[
 \iota_x([U:V])=[U:Uy+V(x-y)].
\]
For $x\ne y$, this identifies $\mathbb P^1$ with the line.
The parameters of $y,x$, and infinity are respectively
$[1:0],[1:1],[0:1]$, independently of $x$.
Every restricted equation is a binary form of degree $L$, whose
coefficients are polynomials of degree at most $L$ in $x$.
The general line is not contained in the base locus. At the chosen
smooth point $y$, outside its other components, such containment
would put its direction in $T_yY$, contradicting the retained
non-tangency condition.

Over $K=\mathbb C(X)$, let $H$ be the greatest common divisor of the
restricted system and put $g=\deg H\le L$.
After division by $H$ the binary linear system has no base point.
Two generic members of it have no common zero: choose one nonzero
member, then choose the second to avoid its finitely many zeros.
Consequently two generic constant complex linear combinations of
the original degree-$L$ forms restrict to binary forms $F_x,G_x$
whose generic gcd is exactly $H$.
The good choices form a nonempty algebraic open set over $K$ and
admit constant complex choices: expanding the coefficients of a
nonzero polynomial over $K$ in a complex-linearly independent basis
shows that it cannot vanish on every tuple of complex constants.
This also covers $g=L$, when the residual forms are constants.

\paragraph{The Sylvester matrix and specialization of its divisor.}
Let $S(x)$ be the $2L$-square matrix of
\[
 K[U,V]_{L-1}\oplus K[U,V]_{L-1}
 \longrightarrow K[U,V]_{2L-1},\qquad
 (a,b)\longmapsto aF_x+bG_x.
\]
Its entries have ambient degree at most $L$ in $x$, and its generic
rank on $X$ is $\rho=2L-g$.
Indeed, after writing $F_x=Hf$ and $G_x=Hh$ with $f,h$ coprime,
the kernel consists of $(hQ,-fQ)$ with $\deg Q=g-1$, and has dimension
$g$ (zero if $g=0$).

Consider an actual center $p$ assigned to $X$ at which the retained
depth over $\overline{\pi_pY}$ exceeds $r$. By the choice of $y$, at least
$r-O(1)$ distinct common roots of $F_p,G_p$ are affine points of
retained components, different from $p,y$, and each lies on no other
component of the fixed base locus. Call these distinguished roots.

The generic gcd divisor specializes with its full length, even when
$X$ is singular at $p$. To see this explicitly, choose an irreducible
algebraic curve through $p$ whose general point lies in the open set
where the preceding generic properties, including the gcd degree,
hold. To obtain it, take successive general linear sections through
$p$ until the local dimension is one, avoiding every
positive-dimensional component of the excluded closed set at each
step. That closed set then has local dimension at most zero.
Choose an irreducible curve component through $p$; it meets the good
open set. Normalize it and use the discrete
valuation ring at a point above $p$. Over its fraction field choose
the gcd form, and scale it so its coefficients are integral and at
least one is a unit. Its reduction $H_0$ is a nonzero homogeneous
form of degree $g$. It divides both specialized forms. In fact, the
quotients are integral: the content law for polynomials over a
discrete valuation ring says that a primitive factor times a quotient
of negative coefficient valuation cannot have all coefficients
integral. Reducing the integral factorizations proves divisibility.
Thus $H_0$ defines a common limiting divisor $D_0$ of length $g$ on
$\mathbb P^1$.

At most $q$ distinguished roots lie in the support of $D_0$.
For this assertion, pass to a finite extension splitting the generic
gcd and a valuation ring above the chosen one. Each distinct generic
root has a unique projective specialization, obtained by scaling its
two coordinates to be integral with one a unit. Each physical root
belongs to some fixed irreducible component of the projective base
locus. Closedness forces its specialization to remain in that
component. A distinguished point lies only on its own retained
component, so a branch specializing to it must already have come
from that retained component. It cannot be either immobile endpoint,
whose fixed parameters specialize to $y$ and $p$, nor a root fixed
at infinity. The generic list bound permits at most $q$ remaining
distinct roots. Their specializations have at most $q$ support points,
regardless of the multiplicities carried by the gcd divisor.

If $F_p,G_p$ are not both zero, their common divisor therefore has
degree at least
\[
 g+r-O(1)-q.
\]
This includes the case that exactly one form is zero: its common
divisor with the other is the other degree-$L$ form, and the matrix
has rank $L$. In all these cases
\[
 \rho-\operatorname{rank}S(p)\ge r-q-O(1)\ge r/2
\]
for large $B$. If both forms vanish, $S(p)=0$ and the rank decrease
is $\rho\ge L\ge r$, giving the same conclusion.
Working throughout with binary forms of degree $L$ records any loss
of an affine leading coefficient as a root at infinity; hence that
phenomenon introduces no exception to the comparison.

We have shown that depth exceeding $r$ forces a matrix rank decrease
of at least $r/2$.  It remains to turn this into a proper polynomial
cut of degree $B\exp(O(s))$, to which center no-capture applies.

\paragraph{Multiplicity compression and the center no-capture bound.}
Choose a $\rho$-square minor $Q(x)$ nonzero generically on $X$.
It has degree at most $2L^2$ and has ambient multiplicity at least
$r/2$ at every excessive actual center. For clarity, at such a center
the selected submatrix has rank at most $\rho-r/2$. Constant
invertible row and column operations put its value into rank normal
form. At least $\lceil r/2\rceil$ rows then have every entry in the ambient
maximal ideal at $p$, so every determinant term belongs to its
$\lceil r/2\rceil$-th power.
No smoothness of $X$, or agreement between ambient generic rank and
generic rank on $X$, is needed for this conclusion.

Choose a test point of $X$ at which $Q$ is nonzero. The multiplicity
compression lemma, Lemma~\ref{poly:compression}, produces a polynomial
which vanishes at all these excessive centers, is nonzero at the
test point, and has degree
\[
 O(L^2/r)=B\exp(O(s)).
\]
It is consequently a proper cut on $X$.
For each fixed $Y$, these cuts can be chosen separately on the center
components $X$ with $Y\in\mathcal Y_X$. Their degree bound is uniform.
The componentwise pooling and no-capture statement on the chosen
center plateau makes the captured center mass $o(|P_x|)$.
This is uniform in the choice of $Y$: a contrary sequence of targets
would contradict the same no-capture statement, which allows arbitrary
polynomials of the specified degree. Averaging with the normalized
weights $n_Y/|P_y|$ therefore removes only $o(1)$ of all incidences.

On the remaining incidences, every group in the retained part of
$\mathcal Y_{X(p)}$ at $p$ has depth at most $r$.
Deleting further targets cannot increase depth.
Apply the row deletion and \eqref{dep:depth-to-degree}. Since rounding
$r$ costs a factor $1+o(1)$, this proves \eqref{dep:wide-R} and
completes the proof of the proposition.
\end{proof}

\begin{remark}\label{dep:uniformity-remark}
The proof bounds incidence mass before discarding centers. It does
not claim a uniform depth estimate for every individual original
target at every center. The conclusion needed later is that, after
one fixed deletion, every target subcollection
of fixed positive row mass satisfies the projected-degree bound.
This permits adaptive choices such as the family on which later
regularity equations fail. It also explains why no bound on the
number of target components enters a union bound.
\end{remark}


\subsection{Rank propagation for separated windows}

Let $v$ be the source width and $s$ the target width, with
$v/s\to0$ or $v/s\to\infty$.  Both are $o(\log B)$.
Their limiting normalized profile $(b_i)$ is the same, but
$T_s,D_i,e_j,\sigma$ below refer only to the target window.
The source and target piece sizes are respectively
$N\exp(-o(v))$ and $N\exp(-o(s))$; we do not replace the
former error by $o(s)$ when $v\gg s$.

Use the retained target set and fixed plateau assignments in
Proposition~\ref{dep:projected-degree}.  Its final one-dimensional
assignment is the persistent global terminal assignment, after the
negligible omission of globally unassigned targets.  Apply the proposition
for the finite batch of plateaux and logarithmic tolerances needed in an
argument, take the union of its negligible incidence deletions, and apply
row Markov once.  Call the remaining source centers typical.  This set is chosen independently
of the later fixed density threshold $\rho>0$; only the large-index
threshold may depend on $\rho$.  Thus the same typical centers serve
every later subcollection of fixed positive target mass, including one
chosen after the center and a testing-degree envelope.  On a plateau of
dimension $l=d-j$, every retained individual carrier projects with
dimension $l$, and such a subcollection has distinct projected degree at least
\begin{equation}\label{sc:projected-degree-interface}
 (e_j/R)\exp(-o(s)).
\end{equation}
The permitted depth denominators are
\begin{align}
 R&=D_1 &&\text{always},\label{sc:R-first}\\
 R&=B\exp\left(-s\sum_{i>j}b_i\right)
       &&\text{if }v/s\to0,\label{sc:R-narrow}\\
 R&=B\exp(-2s)
       &&\text{if }v/s\to\infty,\ b_{j+1}>0,
                           \ l<d-1.\label{sc:R-wide}
\end{align}
Here a row consists of the targets paired with one fixed center $p$.
For a chosen target subset $S$ in that row, $H(T)$ now denotes
$H_{\pi_p(S)}(T)$, the restriction rank of its directions.

\begin{lemma}[Strict rank margin on target plateaux]\label{sc:separated-rank}
At typical centers, ranks of directions to arbitrary dense
remaining target subsets satisfy
\begin{equation}\label{sc:strict-rank}
 H(T_s(b))\ge\frac{\sigma(b)T_s(b)}A\exp(c s)
\end{equation}
with a fixed $c>0$ through all plateaux up to the final
endpoint, except possibly the final one-dimensional
plateau when the center window is narrower.  The constant
may be decreased finitely many times during propagation.
The assertion is needed only at fixed interior testing
times and their one-sided limiting endpoints.
\end{lemma}

\begin{proof}
On the initial interval target no-capture, applied to the
pullback of a direction polynomial, gives rank comparable
to $T^{d-1}$.  Its ratio to $\sigma(b)T/A$ is
\[
 \frac{A}{T^2}=\exp(-2bs).
\]
This has a strict positive margin near the first cut,
since $b_1<0$.

On a plateau of dimension $l$, consider the testing-degree envelope of
the directions to the current row subset.  For each fixed assigned target
carrier whose projection is not contained in this envelope, a generic
linear combination of its pulled-back equations is a proper cut on that
carrier.  The combination can be chosen simultaneously for the finite
carrier list.  Use a slightly later interior time for the fixed pullback
and pooling costs.  Every counted target on such a carrier lies in the
cut.  The individual-cut conclusion
of Proposition~\ref{prof:regularity} therefore removes only $o(|P_y|)$
such targets, uniformly in the center and the envelope.  Hence the
complete projection of the assigned carrier is contained in the envelope
at almost all remaining counted targets.

For this fixed envelope, choose one irreducible component containing each
distinct projected carrier occurring at those targets, and let the targets
inherit the choice for their projected carrier.  Group all envelope
components by dimension, counting targets by these choices; components
receiving no targets remain in the degree of their group.  Some group of
dimension $h\ge l$ carries a fixed positive fraction of the target mass.

If $h>l$, that group borrows the preceding bound with growth exponent
$h\ge l+1$.  This preserves the required growth exponent of $\sigma T$
on the plateau.  Approximate the preceding endpoint closely enough that
the lost logarithmic tolerance is a small fraction of the existing $c$.

If $h=l$, each projected carrier assigned to this group equals its chosen
envelope component, since both are irreducible of dimension $l$.
The projected degree lower bound~\eqref{sc:projected-degree-interface}
and the envelope Hilbert bound give
\[
 H(T)\gtrsim T^l e_j/R\,\exp(-o(s)).
\]
Relative to $\sigma T/A=e_jT^{l+1}/A$, this supplies the
ratio $A/(RT)$, to exponential accuracy.  It suffices to
have a strict negative normalized logarithm for
$R D_{j+1}/A$, the worst point at the next endpoint.
The three choices give, respectively,
\begin{equation}\label{sc:strict-margins}
 b_1+b_{j+1},\qquad
 -\sum_{i>j+1}b_i,\qquad
 b_{j+1}-2.
\end{equation}
The first is strictly negative on the hypersurface
plateau, by $b_1+b_2<0$, and whenever $b_{j+1}\le0$.
For wider centers all remaining cases have
$b_{j+1}>0$ and $l<d-1$, so the third choice applies and
has margin at least one.  For narrower centers and
$b_{j+1}>0$, the second choice has a strict margin when
$l\ge2$, since its sum includes a positive later entry.
Only $l=1$ gives zero rather than a strict margin.

There are at most $d$ distinct cut times.  Choose all
endpoint approximations and fixed logarithmic tolerances
smaller than a sufficiently small fraction of the
finitely many positive margins just displayed.  Decrease
$c$ when necessary.  This completes the propagation.
The degree bounds and prior rank statements quantify
over every dense restriction, so choosing a different
dimension group or dense subset for different centers
does not change the argument.
\end{proof}

For directions covered by $C_p$ the upper bound is
\begin{equation}\label{sc:row-upper}
 H(T)\le CE(T+1).
\end{equation}
If the strict margin reaches the last endpoint, choose a
fixed time just below it with $\sigma(b)/N\ge\exp(-c s/2)$
to exponential accuracy.  Equations~\eqref{sc:strict-rank}
and~\eqref{sc:row-upper} then have a ratio tending to
infinity, at least $\delta^{-1}\exp(c s/3)$ for large
indices.  This contradicts the positive density of
covered rows.

\subsection{The final curve plateau for narrower centers}

At the remaining zero-margin endpoint, all but $o(|P_y|)$ of the covered
targets in each typical row will have their fixed assigned terminal curve
project onto a component of the covering family $C_p$.  These are exactly
the incidences excluded by the global terminal deletion.

The only remaining case has $v/s\to0$ and a final
one-dimensional target plateau, $b_{d-1}<b_d$.  The
preceding plateau has the strict margin in
Lemma~\ref{sc:separated-rank}.  Choose a fixed time $b_0<b_d$ sufficiently
close to $b_d$ within this last plateau, and then choose fixed times
\[
 b_0<b_{\rm eq}<b_{\rm test}<b_d.
\]
Put $\eta=b_{\rm eq}-b_0>0$, $T_0=T_s(b_0)$, and
$T_1=T_s(b_{\rm eq})=T_0\exp(\eta s)$.  Fix $0<\varepsilon<\eta$.
Include this tolerance for the final plateau in the finite depth batch,
and use the typical-center set for the combined batch.  Its deletion is
made before the envelopes below are chosen.  Apply the preceding rank
argument with this combined finite batch; its bounds then hold for every
fixed positive-density restriction.  Only now, for each typical center, let
\[
 V_p=\operatorname{Env}_{T_0}(C_p)
\]
be the envelope of the entire directional family, not
just of its observed incidences.  Its testing rank is
at most $CE(T_0+1)$.

First, a component group of $V_p$ of dimension $h\ge2$
cannot contain covered targets of mass at least $\rho|P_y|$ for any
fixed $\rho>0$.  Indeed take that actual subset and borrow its strict rank
bound from just before the final plateau.  The Hilbert bounds for
the envelope group propagate this with exponent
$h\ge2$, exactly the required exponent of $\sigma T$
on a one-dimensional plateau.  At $b_0$ sufficiently
close to $b_d$ it contradicts~\eqref{sc:row-upper}.
If their target mass were not $o(|P_y|)$ uniformly over the typical
centers, a subsequence would supply one such fixed $\rho$ and contradict
this bound on the same typical-center set.  Thus we may discard all
actual labels on the larger components at a cost of $o(|P_y|)$ per row,
including labels at their intersections with curves;
we do not try to cut those intersections by equations
of possibly excessive degree.

Let $\mathcal F_p$ be the family of every one-dimensional irreducible
component of $V_p$, whether or not it is already represented by an
assigned target projection.  We will impose common equations on this
whole family so that they also contain any covering component through a
surviving label.  The envelope Hilbert lower bound for its entire
one-dimensional component group and the rank upper bound imply
\begin{equation}\label{sc:last-family-degree}
 \deg\mathcal F_p=O(E).
\end{equation}

We next place the projection of the fixed assigned terminal curve in
this family.  For each such carrier $Y$ whose projection is not contained
in $V_p$, a generic linear combination of the pulled-back envelope
equations is a proper cut on $Y$.  Choose the combination simultaneously
for the finite carrier list, and use a slightly later interior time for
the fixed pullback and pooling costs.  Every covered target assigned to
one of these carriers lies in that cut.  The individual-cut conclusion of
Proposition~\ref{prof:regularity} therefore removes only $o(|P_y|)$ such
targets, uniformly in $p$.  This applies to the persistent global
assignment fixed before the depth bounds, including targets at physical
carrier intersections.

At every remaining covered target its assigned projection is contained in $V_p$.
It has dimension one because we retained the depth proof's deletion of
every dimension-dropping incidence.  An irreducible component of $V_p$
containing this projected curve cannot have higher dimension: it would
also contain the observed label, which was excluded from every such
component.  Hence the assigned projection is itself a member of
$\mathcal F_p$.

The components of $\mathcal F_p$ are isolated components of the equations
of degree $T_0$ defining $V_p$, even though that locus
may also have higher-dimensional components elsewhere.
Lemma~\ref{poly:regularity} gives equations of degree at most $T_1$
vanishing on the whole family $\mathcal F_p$, with generic full rank except on
components of total degree
\begin{equation}\label{sc:last-exceptional-degree}
 O(E T_0/T_1)=O(\delta(N/A)\exp(-\eta s)).
\end{equation}
These exceptional components carry negligible target
mass.  This assertion uses the dense-subcollection
uniformity of the depth lemma: on this plateau
\[
 \frac{e_{d-1}}{R}
 =\frac{N/D_d}{B\exp(-b_d s)}=\frac NA.
\]
Thus any subcollection carrying a fixed positive fraction of $|P_y|$
has projected degree at least
$(N/A)\exp(-\varepsilon s)$ for the fixed tolerance selected before
the depth deletion.  Compare
with~\eqref{sc:last-exceptional-degree}.  The ratio is
$O(\delta\exp(- (\eta-\varepsilon)s))\to0$.
Indeed, take the physical assigned carriers whose projections belong to
the exceptional family.  Their distinct projected degree is bounded by
that family's degree, so they cannot carry a fixed positive fraction of $|P_y|$.
If their mass were not $o(|P_y|)$ uniformly over typical centers, one fixed
positive density along a subsequence would contradict this comparison.
This applies to the family selected by the envelope and its equations,
without changing the fixed target assignment or the earlier depth deletion.

The identity $RD_d=A$ explains the room in this comparison: the projection
depth and the last physical cut together cost $A$, while regularity supplies
the strict gap $\eta-\varepsilon>0$ in the logarithmic degree scale.
One must not instead compare
$(N/A)\exp(-o(s))$ directly with $\delta N/A$; the latter
comparison would be unjustified for arbitrarily slowly
vanishing $\delta$.

The regularity lemma supplies one polynomial $g_p$ of degree $O(T_1)$,
a linear combination of maximal Jacobian minors, which is nonzero on
every retained projected component.  Its pullback to each corresponding
fixed assigned physical curve is therefore a proper rational test, with
cleared numerator of degree $O(T_1)$.  Apply the individual-cut conclusion
of Proposition~\ref{prof:regularity}, using the reserved time
$b_{\rm test}$, to remove the zeros of these numerators at $o(|P_y|)$
targets.  This is uniform in the center.  At every surviving label
$g_p\ne0$, so the common equations have full rank there.

At a surviving covered label, take a component $\gamma$
of $C_p$ through it.  Since $C_p\subset V_p$, this curve
lies in a component of $V_p$.  That component cannot have
larger dimension, by the preceding exclusion.  Hence
$\gamma$ is a member of $\mathcal F_p$.  The common
equations have full rank at the label, so this member
equals the projected member of $\mathcal F_p$ through the same label.
That member is the projection of the persistent assigned terminal curve.
Thus this curve projects onto $\gamma\subset C_p$, exactly the terminal incidence
removed in Proposition~\ref{prof:pair-selection}, a
contradiction.

\begin{proof}[Proof of Theorem~\ref{sc:theorem}]
The profile-pair selection exhausts the possibilities.
A positive core is impossible by its isolation degree.
The common-window rank bound contradicts coverage except
at the outermost endpoint, where the lift argument gives
the contradiction.  For separated windows, strict row
rank propagation contradicts coverage except on the
final narrow-center curve plateau, which the preceding
argument resolves by common regularity and terminal
deletion.  Thus the proposed almost-complete sparse-cone
coverage cannot exist.
\end{proof}
